% 第6回　産学連携キックオフ（記入例・模範解答）
\session{6}{2}{11月4日（水）}{産学連携キックオフ}{AIとペルソナ設定}{yes}

\goals{
  \item 企業の課題を、背景、制約、期待されている成果に分けて整理できる
  \item ペルソナとは何か、なぜ設計に必要かを説明できる
  \item AIで複数のペルソナ案を作り、現実味と課題との関連で取捨選択できる
  \item 企業から提供された情報の取り扱いを理解する
}

\begin{caution}{企業の情報の取り扱い}
非公開の資料やデータは、AIに入力しない・SNSや外部に書かない・グループ外に共有しない。判断に迷う場合は担当教員に確認する。\quad \achk 確認した
\end{caution}

\work{事前準備}{課題提示で聞きたいこと（最低1つ）}
\alines{2}{・今の店舗に来店する客の年代構成と、来店から商談に進む割合はどのくらいか。\par
・改装中も営業を続けるのか。提案で変えてよい範囲（建物・外構・運用）はどこまでか。}

\work[80mm]{ワーク1}{課題を整理する}
\instr{\textcolor{ans}{以下は記入例のための架空の課題設定（郊外型ディーラーＡ社の店舗リニューアル）}}
\begin{wstab}{colspec={Q[l,wd=30mm,m] X[l,m]},row{1-Z}={ht=20mm}}
  \hc{背景\newline なぜこの課題か} & \af{来店客の高齢化が進み、20〜30代の来店が少ない。オンラインで下調べを済ませた客が増え、店で何を提供するかを見直したい。} \\
  \hc{制約\newline 予算・敷地・\newline ブランドなど} & \af{既存の建物を活かした改装（建て替えはしない）。改装中も営業を続ける。メーカーのブランド基準（外観・ロゴ）は守る。予算は社内で決まっている（詳細は非公開）。} \\
  \hc{期待される成果} & \af{若い世代や子育て世帯が「行ってみたい」と思える店舗の提案。来店から相談・商談につながる空間と運用の工夫。} \\
  \hc{わからない点\newline 確認したいこと} & \af{現在の来店客の動き（どこで滞在し、どこで帰るか）／スタッフの人数と役割／整備で来る客と購入検討客の割合} \\
\end{wstab}

\work[50mm]{グループ}{メンバーと役割分担}
\begin{wstabh}{colspec={X[2,l,m] X[2,l,m] X[3,l,m]},row{2-Z}={ht=8mm}}
  氏名 & 役割 & 連絡方法・メモ \\
  \af{（自分）} & \af{リーダー・AI活用記録のとりまとめ} & \af{Teams のグループチャット} \\
  \af{メンバーＢ} & \af{調査（市場・競合）} & \af{共有フォルダに資料を集約} \\
  \af{メンバーＣ} & \af{空間・図の担当} & \af{手描きスケッチを写真で共有} \\
  \af{メンバーＤ} & \af{スライド・発表の構成} & \af{} \\
  & & \\
\end{wstabh}

\work[50mm]{ワーク2}{AIでペルソナを作り、2〜3体に絞る}
\begin{promptbox}[ペルソナの案を出す]
自動車ディーラーの店舗設計を提案するために、来店客のペルソナを3種類作ってください。各ペルソナには、年齢、職業、家族構成、居住地の特徴、来店の目的、車選びで重視する点、来店時に感じる不安や面倒、休日の過ごし方を含めてください。互いに大きく異なる3人にしてください。
\end{promptbox}
\begin{promptbox}[ペルソナの現実味を確かめる]
次のペルソナについて、現実の顧客として不自然な点、ステレオタイプになっている点を指摘してください。また、このペルソナが店舗に求めることを、本人の言葉で3つ書いてください。（ペルソナを貼る）
\end{promptbox}
\begin{wstabh}{colspec={X[2,l,m] X[3,l,m] Q[c,wd=14mm,m]},row{2-Z}={ht=10mm}}
  ペルソナ案（一言で） & 不自然な点・ステレオタイプ・課題との関連 & 採用\newline ○× \\
  \af{共働きで子ども1人、初めての買い替え} & \af{課題の中心（子育て世帯）に合う。「母親が子どもの世話」と決めつけた描写は修正} & \ans{○} \\
  \af{社会人2年目、初めて車を買う} & \af{課題の中心（20代）。前半の自分の問いとも重なる} & \ans{○} \\
  \af{定年退職後、趣味でスポーツカー} & \af{現実的だが、課題の対象（若い世代）から外れる} & \ans{×} \\
  \af{車に興味のない大学生} & \af{購入の予定がなく、店舗設計の対象としては弱い} & \ans{×} \\
  \af{地方から転勤してきた会社員} & \af{予備。転勤者は地域の店を知らないという点は活かす} & \ans{△} \\
\end{wstabh}

\work[100mm]{確定}{ペルソナカード}
\newcommand\personacardA[9]{%
\begin{wstab}{colspec={Q[l,wd=26mm,m] X[l,m]},rows={ht=7.5mm}}
  \SetCell[c=2]{bg=phasetint,font=\sffamily\footnotesize\bfseries}#1 & \\
  年齢・職業 & \af{#2} \\ 家族構成 & \af{#3} \\ 居住地の特徴 & \af{#4} \\ 来店の目的 & \af{#5} \\ 重視する点 & \af{#6} \\
  不安や面倒 & \af{#7} \\ 休日の過ごし方 & \af{#8} \\
  \SetRow{ht=14mm}本人の言葉\newline\scriptsize 店舗に求めること & \af{#9} \\
\end{wstab}}
\noindent
\begin{minipage}[t]{0.49\linewidth}\personacardA{ペルソナA　名前：\ans{森 由佳（32）}}
  {32歳・事務職（共働き）}{夫と2歳の子ども}{名古屋市近郊の住宅地。駅から遠く、車が日常の足}
  {軽自動車からミニバンへの買い替え}{安全性、乗り降りのしやすさ、維持費}
  {子どもがぐずると話を聞けない。長く引き止められるのが怖い}{公園、ショッピングモール}
  {「子どもを見ながら、自分のペースで比べたい」\newline「聞きたいときだけ聞ける人がいてほしい」}\end{minipage}\hfill
\begin{minipage}[t]{0.49\linewidth}\personacardA{ペルソナB　名前：\ans{石川 拓海（24）}}
  {24歳・メーカーの技術職（社会人2年目）}{一人暮らし}{郊外の工場に勤務。通勤に車が必要}
  {初めての車の購入。候補を2台に絞って試乗}{価格、燃費、デザイン}
  {営業トークが苦手。値段の相場がわからず不安}{動画で車のレビューを見る、ドライブ}
  {「調べたことを確かめに行きたい」\newline「買うと決めるまで、そっとしておいてほしい」}\end{minipage}\par

\begin{nexttime}
  \item グループでペルソナ2〜3体を確定し、1枚にまとめる
  \item 課題に関して調べておきたいこと（市場、地域、競合）を分担して集める
  \item AI活用記録を、グループの作業でも続ける方法を決める
\end{nexttime}
\begin{wstabh}{colspec={Q[l,wd=20mm,m] X[3,l,m] X[1,l,m]},row{2-Z}={ht=9mm}}
  分野 & 調べること & 担当 \\
  市場 & \af{若年層・子育て世帯の車の購入動向（公的調査・業界団体の調査）} & \af{メンバーＢ} \\
  地域 & \af{店舗周辺の人口構成、住宅地の分布、公共交通の便} & \af{自分} \\
  競合 & \af{近隣の他社ディーラー、異業種の「居心地のよい店」の事例} & \af{メンバーＣ・Ｄ} \\
\end{wstabh}
\aimemoA{ペルソナ案の作成と、ステレオタイプの点検}
{「互いに大きく異なる3人のペルソナを」「不自然な点、ステレオタイプになっている点を指摘して」}
{5案から2案を採用。課題の対象から外れる案は不採用。「母親だけが子どもの世話をする」描写は、AIの指摘も受けて修正した。}
{ペルソナの設定（通勤に車が必要など）が地域の実情に合うかを、市の統計とメンバーの実感で確かめた。企業の非公開資料は入力していない}

\begin{point}
  \item 記入例の課題（Ａ社の店舗リニューアル）は\textbf{架空}。実際の課題に合わせて読み替える。
  \item 課題の整理は「背景・制約・期待される成果」に分けられているかを見る。制約（変えられないもの）を書き落とすと、後で実現性を問われる。
  \item ペルソナは、課題の対象に合っているか、ステレオタイプになっていないかの2点で取捨選択できているかを評価する。不採用の理由が書けていることが大切。
  \item 「本人の言葉」は、後のコンセプトや動線の検証で何度も使う。具体的な一言になっているかを見る。
  \item 企業の情報の取り扱いの確認（チェック欄）を、全員が済ませたかを授業内で確かめる。
\end{point}
